Hilbert（希尔伯特）的计划旨在表明，所有数学都可以在一种形式语言（例如算术语言或集合论语言）中被形式化为一个公理化理论。他相信这样的理论将是\textbf{完备的}。也就是说，对每个语句~$!A$，要么 $\Th{T} \Proves !A$，要么 $\Th{T} \Proves \lnot !A$。在这个意义上，$\Th{T}$ 就能解决每一个数学问题：它要么证明其为真，要么证明其为假。如果 Hilbert 是对的，那么数学也将是\textbf{可判定的}。这是因为任何可公理化的理论都是\textbf{可计算枚举的}，即存在一个可计算函数，能列出该理论的所有定理。我们可以通过逐一列出所有定理来检验一个语句~$!A$ 是否为定理：一直列举，直到找到~$!A$（此时它是定理）或 $\lnot !A$（此时它不是定理）。遗憾的是，Hilbert 错了。Gödel证明了，任何“足够强”的一致且可公理化的理论都不是完备的。这就是\textbf{第一不完备性定理}。“足够强”的要求实质上在于该理论能表示所有可计算函数和关系。具体来说，非常弱的理论 $\Th{Q}$ 就满足这一性质，任何至少和~$\Th{Q}$ 一样强的理论也同样满足。他还证明了——这就是\textbf{第二不完备性定理}——表达该理论一致性的语句在该理论中本身就是不可判定的，即理论既不能证明它，也不能证明其否定。因此，Hilbert 的进一步目标——为整个数学找到“有穷的”一致性证明——也无法实现。因为任何有穷的一致性证明想必都可以在一个涵盖所有数学的理论中被形式化。最后，我们已经证明，能表示所有可计算函数和关系的理论都不是\textbf{可判定的}。需要注意的是，虽然可公理化性和完备性合起来蕴涵可判定性，但不完备性本身并不蕴涵不可判定性。因此，这一结果表明 Hilbert 的第二个目标——即存在一个程序，能够判定是否 $\Th{T} \Proves !A$——也无法实现，至少对于那些不弱于~$\Th{Q}$ 的理论而言是如此。